\section{The polar Lie algebra}
\label{sec:polar}

\begin{proposition}\label{prop:cycles}
Every solution of \eqref{eq:functional} satisfies
\[
 f^3=\id,\qquad f(0)=0,\qquad
 f(x+f(x))=x.
\]
Its only fixed point is zero. Every nonzero orbit spans a binary plane,
and these planes partition \(V\setminus\{0\}\). In particular \(\dim V\)
is even.
\end{proposition}
\begin{proof}
Evaluate \eqref{eq:functional} at \(f(x)\). It gives
\(f^3(x)=f^2(x)+f(x)=x\), so \(f\) is bijective.
If \(a=f(0)\), then \(f(a)=a\); applying \(f^3\) at zero gives \(a=0\).
A fixed point in \eqref{eq:functional} must be zero.
The remaining identity follows from
\(x+f(x)=f^2(x)\). The nonzero orbits have length three and sum zero,
and \(3\mid 2^{\dim V}-1\) implies the parity assertion.
\end{proof}

\begin{theorem}\label{thm:polar}
Let \(f\) be a quadratic solution. Its polar form
\([x,y]=f(x+y)+f(x)+f(y)\) is a Lie bracket. It satisfies
\begin{align}
 [x,f(x)]&=0,\label{eq:commuting}\\
 f([x,y])&=[f(x),f(y)]+[x,y]
                  +[f(x)+f(y),[x,y]],\label{eq:bracket-image}\\
 [[a,b],[c,d]]+[[a,c],[b,d]]+[[a,d],[b,c]]&=0.
 \label{eq:four-variable}
\end{align}
\end{theorem}
\begin{proof}
The polar is alternating and bilinear because \(f\) is quadratic and
\(f(0)=0\). Expand \(f(x+f(x))=x\) by polarization. The terms
\(f(x)+f^2(x)\) sum to \(x\), leaving \eqref{eq:commuting}.
Take the third additive difference at zero of \(x\mapsto[x,f(x)]\).
The terms involving individual values of \(f\) cancel, leaving
\[
 [x,[y,z]]+[y,[z,x]]+[z,[x,y]]=0.
\]
This is Jacobi.

Next polarize \(f^2=f+\id\). Since
\(f(x+y)=f(x)+f(y)+[x,y]\), the expansion of the left side gives
\eqref{eq:bracket-image}.

For \eqref{eq:four-variable}, introduce Boolean variables
\(t_1,\ldots,t_4\) and vectors \(x_1,\ldots,x_4\). The exact expansion is
\[
 f\Bigl(\sum_i t_ix_i\Bigr)
   =\sum_i t_if(x_i)+\sum_{i<j}t_it_j[x_i,x_j].
\]
In applying \(f\) again, the coefficient of \(t_1t_2t_3t_4\) can only
come from a polar cross term between two disjoint pair summands.
There are three such pairings, giving the left side of
\eqref{eq:four-variable}. The composition has degree at most two,
so this coefficient vanishes. All coefficient comparisons take place in
the Boolean function algebra, with \(t_i^2=t_i\).
\end{proof}

The two polarization steps have distinct roles. Binary degree two gives
an alternating bilinear form; the commuting identity
\eqref{eq:commuting} forces its Jacobi identity. The fourth polarization
then supplies \eqref{eq:four-variable}, an additional restriction on the
resulting Lie algebra.
The polynomial in \eqref{eq:four-variable} is the associative standard
polynomial of degree four expressed as a Lie element in characteristic
two~\cite[Theorem~9]{Drensky2021}. Its forced occurrence here gives the
explicit exclusions in Section~\ref{sec:classification}.

\begin{proposition}\label{prop:invariant-series}
All terms of the derived, lower central, and upper central series of the
polar Lie algebra \(L=(V,[\ ,\ ])\) are \(f\)-invariant.
The center is \(f\)-invariant and \(f\) is linear on it.
Every nonzero factor of any of these series has even binary dimension.
\end{proposition}
\begin{proof}
Let \(I\) be an \(f\)-invariant ideal. In
\eqref{eq:bracket-image}, take \(x\in I\). Each term on the right belongs
to \([I,L]\), which is an ideal. Quadratic expansion extends invariance
from bracket generators to their span: the extra polar terms also lie in
that ideal. The same argument with \(x,y\in I\) proves invariance of
\([I,I]\), using that it is an ideal of \(L\).
Induction starts with \(I=L\).

Polarization of \eqref{eq:commuting} gives
\begin{equation}\label{eq:commuting-polarized}
 [x,f(y)]+[y,f(x)]+[x,[x,y]]+[y,[x,y]]=0.
\end{equation}
If \(x\in Z(L)\), this says \([y,f(x)]=0\) for every \(y\).
Thus \(Z(L)\) is invariant; its polar is zero, so \(f\) is linear there.
Pass successively to quotients by invariant central ideals to obtain
the upper central assertion. Every series factor inherits a solution of
\eqref{eq:functional}, so Proposition~\ref{prop:cycles} proves its parity.
\end{proof}

\begin{lemma}[Ideal descent]\label{lem:ideal-descent}
For every Lie ideal \(I\) of \(L\),
\[
 \{f(x)+I:x\in I\}\subseteq Z(L/I).
\]
Every perfect ideal is \(f\)-invariant. Every ideal \(I\) for which
\(Z(L/I)=0\) is also \(f\)-invariant.
\end{lemma}
\begin{proof}
For \(x\in I\), all terms in \eqref{eq:commuting-polarized} except
\([y,f(x)]\) belong to \(I\). This proves the inclusion.
The map \(\phi:I\to L/I\), \(x\mapsto f(x)+I\), is binary-linear
because its polar takes values in \(I\).
Equation~\eqref{eq:bracket-image} gives
\(\phi([x,y])=[\phi(x),\phi(y)]\) for \(x,y\in I\).
Its image is central, so \(\phi([I,I])=0\). This proves the perfect-ideal
assertion. The centerless-quotient assertion follows immediately from
the inclusion. Bijectivity and finiteness turn each invariant inclusion
into equality.
\end{proof}

\begin{corollary}\label{cor:simple-reduction}
If a quadratic solution has a nonsolvable polar algebra, then there is a
quadratic solution whose polar algebra is nonabelian simple, of even
binary dimension no larger than the original dimension.
\end{corollary}
\begin{proof}
The stabilized derived ideal \(P\) is nonzero and perfect, hence invariant.
Choose a maximal proper Lie ideal \(M\) of \(P\). The quotient \(P/M\)
is nonabelian simple and centerless. Apply Lemma~\ref{lem:ideal-descent}
inside \(P\) to see that \(M\) is invariant. Restriction and quotient
therefore give the required quadratic solution.
\end{proof}

\begin{lemma}\label{lem:two-generators}
A nonabelian nilpotent polar Lie algebra has binary abelianization
dimension at least four.
\end{lemma}
\begin{proof}
The abelianization is a nonzero even-dimensional space with the linear
action induced by \(f\). If its dimension were two, any nonzero coset
\(x+L'\) and its image \(f(x)+L'\) would form a binary basis.
The subalgebra \(H\) generated by \(x,f(x)\) is abelian by
\eqref{eq:commuting}. Moreover \(L=H+L'\).
Inductively, \(L=H+\gamma_j(L)\) implies
\(L=H+\gamma_{j+1}(L)\), because
\([H+\gamma_j,H+\gamma_j]\subseteq H+\gamma_{j+1}\).
Nilpotency then gives \(L=H\), a contradiction.
\end{proof}

\begin{example}\label{ex:regular-cpp}
Let \(\F_{16}\) contain \(\F_4\), let \(\eta\in\F_4\setminus\F_2\),
and write \(\chi(x)=1\) off \(\F_4\) and \(\chi(x)=0\) on \(\F_4\).
The affine-switching construction of
\cite[Corollary~6]{Xu2022} includes
\[
 G(x)=\eta x+\chi(x)
      =\eta x+\operatorname{Tr}_{\F_{16}/\F_4}(x)^3.
\]
The trace is binary-linear, and its cube is binary-quadratic.
Both \(\F_4\) and its complement are invariant under multiplication
by \(\eta\) and under translation by one. Hence \(G\) is a permutation
with \(G^3=\id\) and only zero fixed. The same holds for \(G+\id\).
Nevertheless
\[
 G^2(x)+G(x)+x=\eta\chi(x),
\]
so \(G\) does not satisfy \eqref{eq:functional}.
This separates the equation studied here from simultaneous triple-cycle
complete permutation behavior even within binary degree two.
\end{example}
